% Section 6 translation from the user-provided paper.
\clearpage
\section{结论}\label{sec:conclusion}

本文提出了一种新的灵活技术，可以方便地完成相机标定。该技术只要求相机从少数几个（至少两个）不同方向观察一个平面模板。相机或平面模板均可移动，且不需要知道其运动。方法对径向镜头畸变进行了建模。所提出的流程由闭式解和随后基于最大似然准则的非线性细化组成。

我们使用计算机模拟和真实数据对所提出的技术进行了测试，并取得了非常好的结果。与需要使用两个或三个相互正交平面等昂贵设备的经典技术相比，所提出的技术具有显著的灵活性优势。

\section*{致谢}\label{sec:acknowledgment}

感谢 Brian Guenter 提供角点提取软件并进行了多次讨论，感谢 Bill Triggs 提供富有启发性的评论。感谢 Andrew Zisserman 使我注意到他的 CVPR98 工作~\cite{liebowitz1998metric}；该工作使用了相同的约束，但形式不同。感谢 Bill Triggs 和 Gideon Stein 建议开展关于模型不精确性的实验，该实验见技术报告~\cite{zhang1998flexible}。Anandan 和 Charles Loop 审阅了早期版本的英文。
